% Generated draft: requires mathematical and source review.
\documentclass[11pt]{article}
\usepackage[margin=25mm]{geometry}
\usepackage{fontspec}
\setmainfont{Latin Modern Roman}
\usepackage{amsmath,amssymb,mathtools,unicode-math}
\setmathfont{Latin Modern Math}
\usepackage{enumitem,needspace,xurl,hyperref}
\setlength{\parindent}{0pt}
\setlength{\parskip}{5pt}
\setlength{\emergencystretch}{3em}
\newcommand{\N}{\mathbb{N}}
\newcommand{\Z}{\mathbb{Z}}
\newcommand{\Q}{\mathbb{Q}}
\newcommand{\R}{\mathbb{R}}
\newcommand{\C}{\mathbb{C}}
\newcommand{\F}{\mathbb{F}}
\newcommand{\A}{\mathbb{A}}
\newcommand{\PP}{\mathbb{P}}
\newcommand{\E}{\mathbb{E}}
\newcommand{\eps}{\varepsilon}
\newcommand{\dd}{\,\mathrm{d}}
\begin{document}
\section*{UnsolvedMath ID: 800281}
\textbf{Source problem code:} catalogue entry
\subsection{Definitions and mathematical statement}
In the usual several-complex-variables interpretation, the source notation \(P^n\) denotes complex projective space \(\mathbb{P}^n(\C)\); the source does not explicitly specify the underlying field. For a smooth real submanifold \(M\) of this space, let \(J\) be the ambient complex structure and define its complex tangent space by
\[
H_pM=T_pM\cap J(T_pM),\qquad p\in M.
\]
The submanifold is a CR manifold when these spaces have constant complex dimension, called its CR dimension. The source does not specify the differentiability class or the CR dimension of \(M\).

The problem attributed in the source to N. Sibony is the following: let \(M\subset P^n\) be a compact CR manifold of odd real dimension. Find conditions on \(M\) ensuring that \(M\) bounds a complex variety in the ambient projective space. Symbolically, the requested conclusion is
\[
\text{there exists a complex variety }V\text{ with }\partial V=M.
\]
Here \(\partial V=M\) records the source's intended boundary relation; it is not a fully specified boundary operator. The source does not state whether the filling is understood as a variety with a classical boundary, an analytic set away from \(M\) with prescribed limiting boundary, or a holomorphic chain whose boundary is taken in the sense of currents.

Writing \(\dim_{\R}M=2q-1\), the usual pure-dimensional filling would have complex dimension \(q\), but purity, multiplicities, allowable singularities, and the precise boundary relation are not specified. Neither an orientation nor a maximal-complexity hypothesis \(\dim_{\C}H_pM=q-1\) is explicitly supplied, so these are not imposed here. The source also leaves the desired type of conditions unspecified: it asks for conditions ensuring existence, without explicitly demanding a necessary-and-sufficient characterization.

The parenthetical assertion that the analogous question is well understood in \(\C^n\) is preserved as a qualification of the supplied 1995 problem statement, not as an independently verified assessment. A fully precise existence statement requires clarification of the filling and boundary conventions.
\subsection{Short English statement}
N. Sibony asks which conditions ensure that a compact odd-dimensional CR manifold in projective space bounds a complex variety. The source leaves the regularity, CR dimension, orientation, and meaning of the filling and its boundary unspecified. It describes the corresponding question in \(\C^n\) as well understood. The registry labels this problem open; this is supplied status information, not an independently verified mathematical assessment.
\subsection{Sources}
\begin{itemize}
\item Local UnsolvedMath registry, permanent ID \(800281\), source problem code \texttt{catalogue entry}, title ``Several Complex Variables Problem List --- Problem 81.'' The supplied record belongs to the AMR Open Problem Lists collection and carries the status \texttt{open}; its update timestamp is \texttt{2026-10-16T14:00:00Z}. These are supplied registry metadata.
\item Buzzard, Fornaess, and Krantz, \emph{The Several Complex Variables Problem List} (1995), Problem 81. Supplied source URL: \url{https://arxiv.org/abs/math/9509202}. The background records access on \texttt{2026-10-16} and extraction by \texttt{arxiv\_latex}; these are supplied provenance notes, not evidence of independent source consultation here.
\item The registry's \texttt{proposed\_by} field names Buzzard, Fornaess, and Krantz, while the extracted problem itself is signed N. Sibony, with the supplied contact \texttt{sibony@anh.matups.fr}. The problem-level attribution to Sibony is retained separately from the list attribution.
\end{itemize}
\end{document}
